\documentclass[11pt,a4paper]{article}

% ----------- Packages -----------
\usepackage[margin=0.75in]{geometry}
\usepackage{enumitem}
\usepackage{amsmath}
\usepackage{hyperref}
\usepackage{titlesec}
\usepackage{parskip}

\hypersetup{
    colorlinks=true,
    linkcolor=blue,
    urlcolor=blue,
    pdftitle={Joan Alcaide-Nunez CV},
}

% ----------- Formatting -----------
\setlength{\parindent}{0pt}
\titleformat{\section}{\Large\bfseries\uppercase}{}{0pt}{}[\titlerule]
\titlespacing*{\section}{0pt}{12pt}{8pt}

\titleformat{\subsection}{\large\bfseries}{}{0pt}{}
\titlespacing*{\subsection}{0pt}{8pt}{4pt}

% Clean item spacing
\setlist[itemize]{topsep=0pt,itemsep=2pt,parsep=2pt,leftmargin=1.5em}

\begin{document}

% ----------------- NAME + CONTACT -----------------
\begin{center}
    {\LARGE \textbf{Joan Alcaide-Nunez}} \\
    \vspace{4pt}
    \href{mailto:joan.alcaide@campus.lmu.de}{joan.alcaide@campus.lmu.de} \quad | \quad 
    \href{https://joanalnu.github.io}{Website} \quad | \quad 
    \href{https://github.com/joanalnu}{GitHub} \\
    Last updated: April 2026
\end{center}

% ----------------- EDUCATION -----------------
\section*{Education}
\textbf{Ludwig-Maximilians-Universitat (LMU) Munich} \hfill Sept 2025 -- Present \\
B.Sc. in Physics
%\textit{Key Focus:} Transitioning toward Theoretical Physics; establishing foundations in Hamiltonian mechanics and Lorentz Invariance Violation (LIV) probes.

\textbf{German School of Barcelona} \hfill Sept 2017 -- May 2025 \\
Dual High School Diploma (Germany + Spain), GPA: 1.0/1.0

\textbf{Additional Theoretical Training:}
\begin{itemize}
    \item \textit{Gravity and Black Holes} — Perimeter Institute (GoPhysics!), 2025
		\item \textit{Mathematics, Statistics, and Data Science} — Autonomous University of Barcelona (UAB), 2022
		\item \textit{Algorithms and Programming (C++)} - Polytechnic University of Catalonia (UPC), 2022
		\item \text{League of Codes (C++)} - Harbour Space, 2022-2023
\end{itemize}

% ----------------- RESEARCH EXPERIENCE -----------------
\section*{Research Experience}

\textbf{Theoretical Probes of Spacetime (Independent Study)} \hfill \textit{Proposed July -- Sept 2026} \\
\textit{Supervision: Independent Project}
\begin{itemize}
    \item \textbf{Objective:} Deriving Modified Dispersion Relations (MDR) for high-energy photons in non-standard cosmological backgrounds.
    \item \textbf{Formalism:} Applying Hamilton-Jacobi equations to determine cumulative time-of-flight delays ($\Delta t$) for gamma-rays under Quintessence-based FLRW metrics.
    \item \textbf{Innovation:} Investigating the analytical coupling between the dark energy equation of state ($w$) and Planck-scale energy-dependent velocities.
\end{itemize}

\textbf{GRB Cosmography \& Relativistic Foundations} \hfill June -- July 2025 [cite: 8, 10, 87] \\
\textit{Osservatorio Astronomico di Brera (OAB-INAF), Italy} [cite: 9]
\begin{itemize}
    \item \textbf{Mentorship:} Received training on relativistic fireball effects and the theoretical foundations of Gamma-Ray Bursts (GRBs)[cite: 13, 89].
    \item \textbf{Parameterization:} Constrained cosmological density parameters ($\Omega_m, \Omega_\Lambda$) via $\chi^2$ minimization of empirical energy relations[cite: 11, 87].
    \item \textbf{Synthesis:} Integrated GRB results with simulated populations from Einstein Telescope binary neutron star expectations[cite: 12, 88].
    \item \textit{Supervisors: Dr. Giancarlo Ghirlanda, Dr. Om Sharan Salafia} [cite: 14]
\end{itemize}

\textbf{CAPIBARA Collaboration (Founder \& PI)} \hfill 2024 -- Present [cite: 15, 17, 90]
\begin{itemize}
    \item Leading student collaboration in developing instrumentation and mission concepts for high-energy astronomy CubeSat missions[cite: 18, 19, 90].
\end{itemize}

\textbf{Type Ia Supernovae: Mapping the Expansion History} \hfill June -- July 2024 [cite: 25, 27, 93] \\
\textit{Institute of Space Sciences (ICE-CSIC-IEEC), Spain} [cite: 26]
\begin{itemize}
    \item \textbf{Metric Fitting:} Benchmarked the empirical consistency of the local expansion history using SNe Ia light-curve fitting from ESO, ATLAS, and ZTF[cite: 28, 93].
    \item \textbf{Formalism:} Explored statistical foundations of the Hubble-Lemaitre Law through standardizable candle analysis[cite: 29, 94].
    \item \textit{Supervisor: Dr. Lluis Galbany} [cite: 30, 95]
\end{itemize}

\textbf{Cepheid Variable Stars Research} \hfill Aug -- Dec 2023 [cite: 31, 33, 95] \\
\textit{Youth \& Science Program} [cite: 32]
\begin{itemize}
    \item \textbf{Analytical Foundations:} Reconstructed the cosmic distance ladder by deriving the Period-Luminosity relationship from NASA-IPAC NED data[cite: 34, 95].
    \item \textbf{Methodology:} First exposure to rigorous scientific writing and the peer-review process ($H_0 = 70.26 \pm 7.1 \mathrm{\ km\ s^{-1}\ Mpc^{-1}}$)[cite: 34, 35, 96].
\end{itemize}

\textbf{Student Internship at German Space Agency (DLR)} \hfill Oct 2024 [cite: 20, 22, 91] \\
\textit{Institute for Remote Sensing (DLR-IMF), Germany} [cite: 21]
\begin{itemize}
    \item Worked with photogrammetry departments on hyperspectral imaging and LandSat thermal analysis[cite: 23, 91].
\end{itemize}

% ----------------- AWARDS -----------------
\section*{Awards \& Honors}
\textbf{Youth \& Science Fellowship:} 3-year funded fellowship for research stays internationally. \\
\textbf{Abiturpreis German Physical Society (DPG):} Physics Distinction for Academic Achievement in high school. \\
\textbf{IAAC:} Silver Honour and National Award for Spain (2025), top 7\% of >12,000 students. Similar awards in 2023 and 2024. \\
\textbf{Jugend Forscht:} First Prize Iberia (2024/2025), First prizes in regional and national rounds, Special Award in Scientific Photography (NRW 2025).

% ----------------- TALKS -----------------
\section*{Talks \& Outreach}
\textbf{SSEA Symposium:} The CRD Instrument and the COSMOS program, April 2026. \\
\textbf{CosmoXarxa:} High-energy astrophysics presentation at Barcelona Science Museum, March 2025. \\
\textbf{Invited Talk:} Intro to Astrophysics for 11th graders, German School Barcelona, 2024. \\
\textbf{Outreach:} Space Exploration Activity for 3rd grades, Escola Canigó Public School, 2023. \\
\textbf{Explainers (Volunteer):} Science outreach at CosmoCaixa Museum, 2022--2023.
Slides of most talks can be found on my website.

% ----------------- SKILLS -----------------
\section*{Skills}
\textbf{Programming:} Python (\textit{numpy, scipy, pandas, matplotlib, astropy}), C++, LaTeX, Markdown. \\
\textbf{Tools:} Git/GitHub, Visual Studio Code, NeoVim, Linux environment.

\end{document}
